% ═══════════════════════════════════════════════════════════════
% AB-03: Elektrische Leistung & Energie
% Physik Klasse 8 | Greenhouse Schools Rostock
% ═══════════════════════════════════════════════════════════════

\documentclass[11pt, a4paper]{article}

\usepackage[utf8]{inputenc}
\usepackage[T1]{fontenc}
\usepackage[ngerman]{babel}

\usepackage{helvet}
\renewcommand{\familydefault}{\sfdefault}

\usepackage{microtype}
\usepackage[margin=1.5cm, top=2cm, bottom=1.5cm]{geometry}
\usepackage{enumitem}
\usepackage{tabularx}
\usepackage{booktabs}
\usepackage{array}
\usepackage{multicol}
\usepackage{tikz}
\usetikzlibrary{calc}

\usepackage{amsmath, amssymb}
\usepackage{siunitx}
\sisetup{locale = DE, per-mode = fraction, fraction-function = \dfrac}

\usepackage{fancyhdr}
\usepackage{lastpage}
\usepackage{needspace}

\usepackage{xcolor}
\usepackage{tcolorbox}
\tcbuselibrary{skins}

% ═══════════════════════════════════════════════════════════════
% FARBEN
% ═══════════════════════════════════════════════════════════════

\definecolor{physik}{HTML}{2c5aa0}
\definecolor{hellgrau}{HTML}{F5F5F5}
\definecolor{mittelgrau}{HTML}{888888}
\definecolor{rahmen}{HTML}{CCCCCC}

% ═══════════════════════════════════════════════════════════════
% VARIABLEN
% ═══════════════════════════════════════════════════════════════

\newcommand{\fach}{Physik}
\newcommand{\thema}{Elektrische Leistung \& Energie}
\newcommand{\klasse}{8}
\newcommand{\stunde}{03}
\newcommand{\maxpunkte}{25}

\newif\ifzeigesternchen
\zeigesternchentrue

\colorlet{fachfarbe}{physik}

% ═══════════════════════════════════════════════════════════════
% KOPF- UND FUSSZEILE
% ═══════════════════════════════════════════════════════════════

\pagestyle{fancy}
\fancyhf{}
\renewcommand{\headrulewidth}{0.4pt}
\renewcommand{\footrulewidth}{0pt}
\lhead{\small \fach{} Klasse \klasse{} | \thema}
\rhead{\small Name: \underline{\hspace{3.5cm}}}
\cfoot{\small\textcolor{mittelgrau}{Seite \thepage{} von \pageref{LastPage}}}

% ═══════════════════════════════════════════════════════════════
% BOXEN
% ═══════════════════════════════════════════════════════════════

\newtcolorbox{merkkasten}[1][]{
    colback=hellgrau,
    colframe=hellgrau,
    boxrule=0pt,
    arc=0pt,
    left=8pt, right=8pt, top=8pt, bottom=6pt,
    borderline north={2pt}{0pt}{fachfarbe},
    before skip=6pt, after skip=6pt,
    #1
}

\newtcolorbox{formelkasten}{
    colback=white,
    colframe=fachfarbe,
    boxrule=1pt,
    arc=0pt,
    left=8pt, right=8pt, top=6pt, bottom=6pt,
    before skip=4pt, after skip=4pt
}

% ═══════════════════════════════════════════════════════════════
% BEFEHLE
% ═══════════════════════════════════════════════════════════════

\newcommand{\luecke}[1]{\underline{\hspace{#1}}}
\newcommand{\punktebox}[1]{\hfill\small\textcolor{mittelgrau}{[\_\_/#1]}}
\newcommand{\stern}[1]{\ifzeigesternchen\textsuperscript{\textcolor{fachfarbe}{(#1)}}\fi}

\newcommand{\aufgabe}[2]{%
    \needspace{4\baselineskip}%
    \vspace{10pt}%
    \noindent\textbf{\textcolor{fachfarbe}{Aufgabe #1}} #2%
    \vspace{4pt}%
    \nopagebreak[4]%
}

\newcommand{\operator}[1]{\textbf{\textcolor{fachfarbe}{#1}}}

\newlist{teil}{enumerate}{2}
\setlist[teil,1]{label=\alph*), leftmargin=1.8em, itemsep=3pt, parsep=0pt, topsep=3pt}

\newcommand{\antwortzeilen}[1]{%
    \par\vspace{4pt}%
    \foreach \n in {1,...,#1}{%
        \noindent\rule{\linewidth}{0.3pt}\vspace{6pt}%
    }%
}

\newcommand{\ein}[1]{\,\mathrm{#1}}

\setlength{\parindent}{0pt}
\setlength{\parskip}{4pt}

% ═══════════════════════════════════════════════════════════════
% DOKUMENT
% ═══════════════════════════════════════════════════════════════

\begin{document}

% HEADER
\begin{center}
    {\Large\bfseries\textcolor{fachfarbe}{\thema}}\\[0.2em]
    {\small Arbeitsblatt | \fach{} Klasse \klasse{} | Stunde \stunde}
\end{center}
\vspace{-0.3em}
\hrule
\vspace{0.6em}

% ═══════════════════════════════════════════════════════════════
% MERKKASTEN 1: Leistung
% ═══════════════════════════════════════════════════════════════

\begin{merkkasten}
\textbf{Elektrische Leistung}

Die elektrische \luecke{2.5cm} gibt an, wie viel \luecke{2.5cm} pro \luecke{2cm} umgewandelt wird.

\vspace{0.3em}
\textbf{Definition:} \quad $P = \dfrac{E}{t}$ \quad (Leistung = Energie geteilt durch Zeit)

\vspace{0.3em}
\textbf{Berechnung:} \quad $P = $ \luecke{2cm} \qquad
\textbf{Einheit:} \quad $[P] = 1$ \luecke{1.5cm} $= 1\ein{W}$

\vspace{0.3em}
\textbf{Umrechnung:} \quad $1\ein{kW} = $ \luecke{1.5cm} $\ein{W}$
\end{merkkasten}

% ═══════════════════════════════════════════════════════════════
% MERKKASTEN 2: Energie
% ═══════════════════════════════════════════════════════════════

\begin{merkkasten}
\textbf{Elektrische Energie}

Elektrische Energie ist \luecke{2.5cm} mal \luecke{2cm}. \quad (Umstellung: $E = P \cdot t$)

\vspace{0.3em}
\textbf{Formel:} \quad $E = $ \luecke{2cm}

\vspace{0.3em}
\textbf{SI-Einheit:} \quad $1\ein{Ws} = 1\ein{J}$ \quad (Joule ist die SI-Einheit der Energie)

\vspace{0.3em}
\textbf{Haushaltsüblich:} \quad $1\ein{kWh} = $ \luecke{2.5cm} $\ein{J}$ \quad (für Stromabrechnung)

\vspace{0.3em}
\textbf{Umrechnung:} \quad $1\ein{kWh} = $ \luecke{1.5cm} $\ein{Wh}$
\end{merkkasten}

% ═══════════════════════════════════════════════════════════════
% AUFGABE 1: Geräteetiketten
% ═══════════════════════════════════════════════════════════════

\aufgabe{1}{\stern{*}} \punktebox{6}

\operator{Suche} im Physikraum drei elektrische Geräte mit Typenschild. \operator{Notiere} die Angaben.

\vspace{0.3em}
\renewcommand{\arraystretch}{1.6}
\begin{tabularx}{\linewidth}{|p{3.5cm}|X|X|X|X|}
\hline
\textbf{Gerät} & \textbf{Spannung (V)} & \textbf{Stromstärke (A)} & \textbf{Leistung (W)} & \textbf{Sonstiges} \\
\hline
1. & & & & \\
\hline
2. & & & & \\
\hline
3. & & & & \\
\hline
\end{tabularx}
\renewcommand{\arraystretch}{1.0}

% ═══════════════════════════════════════════════════════════════
% AUFGABE 2: Berechnen
% ═══════════════════════════════════════════════════════════════

\aufgabe{2}{\stern{*}} \punktebox{4}

\operator{Berechne} die elektrische Leistung.

\begin{teil}
    \item Ein Gerät wird mit $U = 12\ein{V}$ und $I = 3\ein{A}$ betrieben.

    \vspace{1.5em}
    \hfill $P = $ \luecke{2cm}

    \item Ein Föhn zieht bei $U = 230\ein{V}$ einen Strom von $I = 10\ein{A}$.

    \vspace{1.5em}
    \hfill $P = $ \luecke{2cm}
\end{teil}

% ═══════════════════════════════════════════════════════════════
% AUFGABE 3: Energie und kWh
% ═══════════════════════════════════════════════════════════════

\aufgabe{3}{\stern{**}} \punktebox{6}

\operator{Berechne} die elektrische Energie in Wh und kWh.

\begin{teil}
    \item Eine Lampe mit $P = 60\ein{W}$ leuchtet $t = 5\ein{h}$.

    \vspace{2em}
    $E = $ \luecke{3cm} $= $ \luecke{2cm} Wh $= $ \luecke{2cm} kWh

    \item Ein Computer mit $P = 200\ein{W}$ läuft $t = 4\ein{h}$ pro Tag, 30 Tage im Monat.

    \vspace{2.5em}
    $E_{\text{Monat}} = $ \luecke{6cm} $= $ \luecke{2cm} kWh
\end{teil}

% ═══════════════════════════════════════════════════════════════
% AUFGABE 4: Stromkosten
% ═══════════════════════════════════════════════════════════════

\aufgabe{4}{\stern{**}} \punktebox{5}

\operator{Berechne} die Stromkosten. Der Strompreis beträgt $\mathbf{30\,\dfrac{ct}{kWh}}$.

\vspace{0.3em}
Ein Durchlauferhitzer hat $P = 20\ein{kW}$. Familie Müller duscht täglich insgesamt 15 Minuten.

\begin{teil}
    \item Wie viel Energie wird pro Tag verbraucht?

    \vspace{2em}
    $E = $ \luecke{6cm} $= $ \luecke{2cm} kWh

    \item Wie viel kostet das Duschen im Monat (30 Tage)?

    \vspace{2.5em}
    Kosten $= $ \luecke{6cm} $= $ \luecke{2cm} Euro
\end{teil}

% ═══════════════════════════════════════════════════════════════
% AUFGABE 5: Transfer
% ═══════════════════════════════════════════════════════════════

\aufgabe{5}{\stern{***}} \punktebox{4}

\operator{Erkläre}, warum Energieversorger die Einheit \textbf{kWh} statt \textbf{Joule} verwenden.

\antwortzeilen{3}

% ═══════════════════════════════════════════════════════════════
% ERWEITERUNG: Haushaltsanalyse
% ═══════════════════════════════════════════════════════════════

\vspace{1em}
\hrule
\vspace{0.8em}

\noindent\textbf{\textcolor{fachfarbe}{Erweiterung (Hausaufgabe):}} \hfill \textit{freiwillig}

\vspace{0.3em}
\noindent\operator{Analysiere} den Stromverbrauch in deinem Haushalt. \operator{Recherchiere} die Leistungsaufnahme typischer Geräte und \operator{berechne} den Jahresverbrauch.

\vspace{0.3em}
\small
\textit{\textbf{Hinweis:} Bei Geräten wie Kühlschrank oder Waschmaschine steht auf dem Typenschild die \textbf{maximale} Leistung. Der tatsächliche Verbrauch ist geringer, da diese Geräte nicht dauerhaft mit voller Leistung laufen. Recherchiere typische \textbf{Durchschnittswerte} (z.B. Kühlschrank $\approx$ 100--150 kWh/Jahr, Waschmaschine $\approx$ 0,5--1 kWh pro Waschgang).}
\normalsize

\vspace{0.5em}
\renewcommand{\arraystretch}{1.8}
\begin{tabularx}{\linewidth}{|p{3.2cm}|X|X|X|X|}
\hline
\textbf{Gerät} & \textbf{Leistung} & \textbf{Laufzeit/Tag} & \textbf{Energie/Jahr} & \textbf{Kosten/Jahr} \\
 & (W oder kW) & (h) & (kWh) & (€, bei 30 ct/kWh) \\
\hline
Kühlschrank & & & & \\
\hline
Fernseher & & & & \\
\hline
Computer/Laptop & & & & \\
\hline
Smartphone-Laden & & & & \\
\hline
Licht (gesamt) & & & & \\
\hline
Waschmaschine & & & & \\
\hline
 & & & & \\
\hline
 & & & & \\
\hline
\multicolumn{3}{|r|}{\textbf{Summe:}} & & \\
\hline
\end{tabularx}
\renewcommand{\arraystretch}{1.0}

\vspace{0.5em}
\textbf{a)} Welche drei Geräte verbrauchen in deinem Haushalt am meisten Strom?

\antwortzeilen{2}

\textbf{b)} \operator{Schlage} zwei konkrete Maßnahmen vor, um Strom zu sparen.

\end{document}
